% preamble.tex -- shared macros for the paper (input by main.tex before \begin{document})
\documentclass[11pt,a4paper]{amsart}
\usepackage[utf8]{inputenc}
\usepackage[T1]{fontenc}
\usepackage{amsmath,amssymb,amsthm,mathrsfs}
\usepackage[margin=1in]{geometry}
\usepackage{enumitem}
\usepackage{xcolor}   % needed for the colour expressions in hyperref's options below
\usepackage[colorlinks=true,linkcolor=blue!50!black,citecolor=blue!50!black,urlcolor=blue!50!black]{hyperref}
\usepackage{cleveref}
\crefname{hypothesis}{Hypothesis}{Hypotheses}
\Crefname{hypothesis}{Hypothesis}{Hypotheses}
\crefname{convention}{Convention}{Conventions}
\Crefname{convention}{Convention}{Conventions}

% ---------- theorem environments (numbered within sections) ----------
\theoremstyle{plain}
\newtheorem{theorem}{Theorem}[section]
\newtheorem{proposition}[theorem]{Proposition}
\newtheorem{lemma}[theorem]{Lemma}
\newtheorem{corollary}[theorem]{Corollary}
\theoremstyle{definition}
\newtheorem{definition}[theorem]{Definition}
\newtheorem{hypothesis}[theorem]{Hypothesis}
\newtheorem{convention}[theorem]{Convention}
\theoremstyle{remark}
\newtheorem{remark}[theorem]{Remark}
\newtheorem{example}[theorem]{Example}

% ---------- sets, fields, operators ----------
\newcommand{\R}{\ensuremath{\mathbb R}}
\newcommand{\C}{\ensuremath{\mathbb C}}
\newcommand{\Z}{\ensuremath{\mathbb Z}}
\newcommand{\N}{\ensuremath{\mathbb N}}
\newcommand{\E}{\ensuremath{\mathbb E}}
\newcommand{\cS}{\ensuremath{\mathcal S}}
\newcommand{\cD}{\ensuremath{\mathcal D}}
\newcommand{\cK}{\ensuremath{\mathcal K}}
\newcommand{\cI}{\ensuremath{\mathcal I}}
\newcommand{\cH}{\ensuremath{\mathcal H}}
\newcommand{\cV}{\ensuremath{\mathcal V}}
\newcommand{\cF}{\ensuremath{\mathcal F}}
\newcommand{\cT}{\ensuremath{\mathcal T}}
\newcommand{\cZ}{\ensuremath{\mathcal Z}}
\newcommand{\cR}{\ensuremath{\mathcal R}}
\newcommand{\cM}{\ensuremath{\mathcal M}}
\newcommand{\cB}{\ensuremath{\mathcal B}}
\newcommand{\cA}{\ensuremath{\mathcal A}}
\newcommand{\cW}{\ensuremath{\mathcal W}}
\newcommand{\cG}{\ensuremath{\mathcal G}}
\newcommand{\cQ}{\ensuremath{\mathcal Q}}
\newcommand{\cC}{\ensuremath{\mathcal C}}
\newcommand{\cL}{\ensuremath{\mathcal L}}
\newcommand{\sK}{\ensuremath{\mathsf K}}
\newcommand{\sS}{\ensuremath{\mathsf S}}
\newcommand{\sW}{\ensuremath{\mathsf W}}
\newcommand{\amp}{\ensuremath{\mathsf a}}
\newcommand{\Wc}{\ensuremath{\mathsf W}}
\DeclareMathOperator{\supp}{supp}
\DeclareMathOperator{\dist}{dist}
\DeclareMathOperator{\Var}{Var}
\DeclareMathOperator{\Cov}{Cov}
\DeclareMathOperator{\vertex}{vert}
\DeclareMathOperator{\alive}{alive}
\DeclareMathOperator{\Tr}{Tr}
\DeclareMathOperator{\diag}{diag}
\DeclareMathOperator{\curl}{curl}
\renewcommand{\Re}{\operatorname{Re}}
\renewcommand{\Im}{\operatorname{Im}}
\newcommand{\eps}{\ensuremath{\varepsilon}}
\newcommand{\Lam}{\ensuremath{\Lambda}}
\newcommand{\ND}{\ensuremath{\mathrm{ND}}}
\newcommand{\Oh}{\ensuremath{O_h}}
\newcommand{\Oth}{O(3)}
\newcommand{\Euc}{\ensuremath{\R^3\rtimes O(3)}}
\newcommand{\Cub}{\ensuremath{\R^3\rtimes O_h}}
\newcommand{\dstar}{\ensuremath{d^{*}}}
\newcommand{\ie}{i.e.\ }
\newcommand{\eg}{e.g.\ }
\newcommand{\lmu}{\ensuremath{\mu_0}}            % the fixed photon mass
\newcommand{\mus}{\ensuremath{\mu_*}}            % the reference photon mass of the frozen family
\newcommand{\dm}{\ensuremath{\delta m^2}}        % the mass counterterm

% ---------- lattices, cubes, states ----------
\newcommand{\Lk}{\ensuremath{\Lambda_k}}                 % \eps_k Z^3
\newcommand{\CJ}{\ensuremath{C_J}}                       % the cube [-MJ,MJ)^3
\newcommand{\CK}{\ensuremath{C^{K}_{J}}}                 % King's RG cube (as a state convention)
\newcommand{\XJ}{\ensuremath{X_J}}                       % C_J \cap \Lambda_k
\newcommand{\TJ}{\ensuremath{T_J}}                       % the three top face layers
\newcommand{\NDK}{\ensuremath{\ND^{K}}}                  % King's convention (scalar killed on T_J)
\newcommand{\NDa}{\ensuremath{\ND^{(a)}}}                % comparison convention (scalar alive on X_J)
\newcommand{\avK}[1]{\ensuremath{\langle #1\rangle^{K}_{k,J}}}
\newcommand{\ava}[1]{\ensuremath{\langle #1\rangle^{(a)}_{k,J}}}
\newcommand{\av}[1]{\ensuremath{\langle #1\rangle}}
\newcommand{\ZK}{\ensuremath{Z^{K}}}                     % King's printed state functional
\newcommand{\ZKinf}{\ensuremath{Z^{K}_{\infty}}}
\newcommand{\ZKJ}{\ensuremath{Z^{K}_{J}}}
\newcommand{\Ginf}{\ensuremath{G_\infty}}                % the tube object (Leg W)
\newcommand{\ZLIM}[1]{\ensuremath{Z^{#1}_{\infty}}}      % LIM state of a continuum subsequence
\newcommand{\Dtube}{\ensuremath{\cD_{\sigma_+,H}}}       % OS0's Wick tube
\newcommand{\Wick}[1]{\ensuremath{{:}\phi^2{:}(#1)}}
\newcommand{\Fcurv}[1]{F(#1)}
\newcommand{\norm}[1]{\ensuremath{\lVert #1\rVert}}
\newcommand{\abs}[1]{\ensuremath{\lvert #1\rvert}}
\newcommand{\wnorm}[1]{\ensuremath{\lvert #1\rvert_{w}}}
\newcommand{\fournorm}[1]{\ensuremath{\lvert #1\rvert_{4}}}
\newcommand{\ip}[2]{\ensuremath{\langle #1,#2\rangle}}

% ---------- named hypotheses (always typeset through these macros) ----------
\newcommand{\hyp}[1]{\ensuremath{(\mathrm{#1})}}
\newcommand{\EC}{\ensuremath{(\mathrm{E}_{C})}}
\newcommand{\ECK}{\ensuremath{(\mathrm{E}_{C})^{K}}}
\newcommand{\ECa}{\ensuremath{(\mathrm{E}_{C})^{(a)}}}
\newcommand{\Wone}{\ensuremath{(\mathrm{W1})}}
\newcommand{\Wtwo}{\ensuremath{(\mathrm{W2})}}
\newcommand{\Tall}{\ensuremath{(\mathrm{T})}}
\newcommand{\Tone}{\ensuremath{(\mathrm{T1})}}
\newcommand{\Ttwo}{\ensuremath{(\mathrm{T2})}}
\newcommand{\Tthree}{\ensuremath{(\mathrm{T3})}}
\newcommand{\BALm}{\ensuremath{(\mathrm{BAL}\text{-}m)}}
\newcommand{\RC}{\ensuremath{(\mathrm{R}_{C})}}
\newcommand{\RW}{\ensuremath{(\mathrm{R}_{W})}}
\newcommand{\RCfr}{\ensuremath{(\mathrm{R}_{C})^{\mathrm{fr}}}}
\newcommand{\RWfr}{\ensuremath{(\mathrm{R}_{W})^{\mathrm{fr},(a)}}}
\newcommand{\TWK}{\ensuremath{(\mathrm{TW}^{K})}}        % item 51: the Wick tube on King's cube
\newcommand{\FAK}{\ensuremath{(\mathrm{FA}^{K})}}        % item 53: Form (A) on King's cube (optional extension)
\newcommand{\Emu}{\ensuremath{(\mathrm{E}_{\mu})}}
\newcommand{\EmuK}{\ensuremath{(\mathrm{E}_{\mu})^{K}}}
\newcommand{\EmuI}{\ensuremath{(\mathrm{E}_{\mu})^{\cI}}}
\newcommand{\EmuW}{\ensuremath{(\mathrm{E}_{\mu})^{W}}}
\newcommand{\NEGz}{\ensuremath{(\mathrm{NEG}_{0})}}
\newcommand{\NEGs}[1]{\ensuremath{(\mathrm{NEG}_{0})[#1]}}
\newcommand{\NEGlim}{\ensuremath{(\mathrm{NEG}_{\lim})}}
\newcommand{\CTmu}{\ensuremath{(\mathrm{CT}_{\mu})}}
\newcommand{\CTmufour}{\ensuremath{(\mathrm{CT}_{\mu})^{(4)}}}
\newcommand{\hypAlph}{\ensuremath{(\alpha)}}   % named hypothesis (alpha); not \Alph, which is LaTeX's counter formatter
% \Alph is LaTeX's counter formatter (\Alph{enumi}); the sections also write \Alph for the hypothesis (alpha).
% Dispatch on the next token: a brace means the counter formatter, anything else the hypothesis.
\makeatletter
\makeatother
\newcommand{\CONT}{\ensuremath{(\mathrm{CONT}_{\mu,J})}}
\newcommand{\Kfr}{\ensuremath{(\mathrm{K}^{\mathrm{fr}}_{\cI})}}
\newcommand{\Hrelunif}{\ensuremath{(\mathrm{H0rel})^{\mathrm{unif}}}}
\newcommand{\Kposunif}{\ensuremath{(\mathrm{K}\text{-}\mathrm{pos})^{\mathrm{unif}}_{N}}}
\newcommand{\KfrWunif}{\ensuremath{(\mathrm{K}^{\mathrm{fr}}_{W})^{\mathrm{unif}}}}
\newcommand{\TR}{\ensuremath{(\mathrm{TR})}}
\newcommand{\TRW}{\ensuremath{(\mathrm{TR})^{W}}}
\newcommand{\Bcell}{\ensuremath{(\mathrm{B}_{\mathrm{cell}})}}
\newcommand{\SCH}{\ensuremath{(\mathrm{SCH})}}
\newcommand{\UNIQ}{\ensuremath{(\mathrm{UNIQ})}}
\newcommand{\PBIR}{\ensuremath{(\mathrm{PB}_{\mathrm{IR}})}}
\newcommand{\PBmus}{\ensuremath{(\mathrm{PB}_{\mus})}}
\newcommand{\NTtau}{\ensuremath{(\mathrm{NT}_{\tau})}}
\newcommand{\hypR}{\ensuremath{(\mathrm{R})}}
\newcommand{\hypRK}{\ensuremath{(\mathrm{R})^{K}}}
\newcommand{\CFK}{\ensuremath{(\mathrm{C}\text{-}\mathrm{FK})}}
\newcommand{\Bminus}[1]{\ensuremath{(\mathrm{B}^{-})[#1]}}
\newcommand{\Itwo}{\ensuremath{(\mathrm{I2})}}
\newcommand{\Ithree}{\ensuremath{(\mathrm{I3})}}
\newcommand{\Ifour}{\ensuremath{(\mathrm{I4})}}
\newcommand{\UVS}{\ensuremath{(\mathrm{UVS}^{+}_{\infty})}}

% ---------- Osterwalder--Schrader axiom labels ----------
\newcommand{\OSz}{\ensuremath{\mathrm{OS0}}}
\newcommand{\OSo}{\ensuremath{\mathrm{OS1}}}
\newcommand{\OSt}{\ensuremath{\mathrm{OS2}}}
\newcommand{\OSth}{\ensuremath{\mathrm{OS3}}}
\newcommand{\OSf}{\ensuremath{\mathrm{OS4}}}
\newcommand{\Ezz}{\ensuremath{\mathrm{E0}''}}

% ---------- small conveniences ----------
\newcommand{\defeq}{:=}
\newcommand{\inv}{\ensuremath{^{-1}}}
\newcommand{\half}{\ensuremath{\tfrac12}}
\newcommand{\lmt}[1]{\ensuremath{\lim_{#1}}}
\allowdisplaybreaks
\providecommand{\cP}{\mathcal{P}}

% ---------- visible placeholder for values or lists still changing (draft only) ----------
\newcommand{\moving}[1]{\textcolor{orange!80!black}{[#1]}}
